\begin{proof}
\begin{enumerate}[label=\arabic*.]
\item \textbf{The middle form and its characteristic vector are fixed.}
The integral calculation in \cref{maps:src:F.2.6} gives a free
middle group of rank four.  The characteristic-number calculation
gives $I_{3,1}$, $c_2^2=2$, and $p_1=-2c_2$, since $H^2=0$.
The reduction of $c_2$ is $w_4$.  The Wu formula on this spin
eight-manifold says that $w_4$ is characteristic for the middle
form.  Hence $v=p_1/2=-c_2$ is a characteristic vector with
$v^2=2$.  It is primitive, since the square of any integer
dividing $v$ would divide two.

\item \textbf{Explicit lattice isometries carry $v$ to
$(1,1,1,1)$.}
Testing the characteristic condition on the four standard basis
vectors of $I_{3,1}$ makes every coordinate of $v$ odd.
Independent sign changes are integral isometries.  After those
changes write $v=(a,b,c,d)$ with positive odd coordinates and
\[
a^2+b^2+c^2=d^2+2.
\]
For $d>1$, put $s=a+b+c$.  We have $s>d$, since
$s^2\geq a^2+b^2+c^2=d^2+2$, and
$s\leq\sqrt{3(d^2+2)}<3d$.  The vector
$u=(1,1,1,1)$ has square two.  Its reflection is the explicit
integral isometry
\[
R_u(z)=z-(z,u)u.
\]
It replaces the fourth coordinate by $2d-s$, whose absolute
value is strictly smaller than $d$.  All four new coordinates
remain odd because $s-d$ is even; in particular the new fourth
coordinate is nonzero.  Taking absolute values by sign
isometries and iterating therefore terminates at $d=1$.
The equation then forces $a=b=c=1$.  The product of these
specified reflections and sign changes is the required
isometry of the lattice together with $p_1/2$.

\item \textbf{The connected sum has the same extended normal form.}
Let $M_0=3\mathbb{HP}^2\#(-\mathbb{HP}^2)$ with its indicated
orientation.  Its standard middle generators have form $I_{3,1}$
and $p_1(TM_0)=2\sum_i x_i$; the normalized tautological-bundle
calculation is \cite[Lemma 2.1 and Theorem 2.4, \S2]{ffMandal}.
Thus the preceding isometry matches both the intersection form
and $p_1$.

For completeness these two pieces of information determine the
unstable normal datum used in the classification.  The normal
rank-four bundle of an embedded middle sphere has Euler number
equal to its self-intersection and Pontryagin number equal to
the evaluation of the ambient $p_1$, since $p_1(TS^4)=0$.
In the clutching group
$\pi_3\operatorname{Spin}(4)=\mathbb Z\oplus\mathbb Z$ the two
numbers, in compatible generators, are
\[
e=a-b,\qquad p_1=-2(a+b).
\]
This homomorphism is injective: its zero equations give
$a=b$ and $a=-b$, hence $a=b=0$.  No further normal-bundle
choice remains after the form and $p_1$ are matched.

\item \textbf{The known smooth classification gives an almost
diffeomorphism.}
Wall's classification, in the explicit extended-form formulation
\cite[Definition 3.6, pp.~12--13, and Theorem 3.7(1),
pp.~13--14]{IVVWallForm}, identifies these oriented
three-connected eight-manifolds up to connected sum with an
oriented homotopy eight-sphere.  Its input is exactly the normal
form identified in the preceding step.  Thus the constructed
manifold is oriented diffeomorphic to $M_0\#\Sigma$ for some
$\Sigma\in\Theta_8$.

\item \textbf{The homotopy sphere is absorbed in a displayed
summand.}
Kramer--Stolz's Theorem A
\cite[Theorem A, p.~178]{IVVKramerStolz} applies to
$\mathbb{HP}^2$: it is simply connected, has dimension eight,
and has integral homology $\mathbb Z$ in degrees zero, four
and eight and zero otherwise.  It gives
$\mathbb{HP}^2\#\Sigma\cong\mathbb{HP}^2$.
This diffeomorphism preserves orientation because both sides
have signature one.  Associativity of the oriented connected
sum and this map give the explicit comparison
\[
M_0\#\Sigma
\cong(\mathbb{HP}^2\#\Sigma)\#\mathbb{HP}^2
       \#\mathbb{HP}^2\#(-\mathbb{HP}^2)
\cong M_0.
\]
Composing it with the almost-diffeomorphism proves the claimed
oriented diffeomorphism.
\end{enumerate}
\end{proof}
